% !TEX program = xelatex
% Generated from 毕业论文初稿.md by generate-latex.mjs.
\documentclass[UTF8,12pt,a4paper,oneside,openany]{ctexbook}
\usepackage[margin=2.5cm,headheight=16pt]{geometry}
\usepackage{booktabs,longtable,array,enumitem,xcolor,fancyhdr,tikz,xurl}
\usetikzlibrary{arrows.meta,shapes.geometric,quotes}
\usepackage{hyperref}
\definecolor{Ink}{HTML}{252731}\definecolor{Line}{HTML}{BAA4AD}\definecolor{Paper}{HTML}{FFF8FA}
\hypersetup{colorlinks=true,linkcolor=Ink,citecolor=Ink,urlcolor=Ink,bookmarksnumbered=true,pdftitle={基于 Node.js 与 MySQL 的在线学习平台设计与实现}}
\setlength{\parindent}{2em}\setlength{\parskip}{.3em}\setlength{\emergencystretch}{3em}
\setlist{leftmargin=2em,itemsep=.2em}\renewcommand{\arraystretch}{1.2}
\newcolumntype{P}[1]{>{\raggedright\arraybackslash}p{#1}}
\pagestyle{fancy}\fancyhf{}\fancyhead[L]{\small 知途课堂毕业设计论文初稿}\fancyhead[R]{\small 2026年10月8日}\fancyfoot[C]{\thepage}
\begin{document}
\begin{titlepage}\centering\vspace*{2cm}
{\LARGE 本科毕业设计论文（初稿）\par}\vspace{1.6cm}
{\Huge\bfseries 基于 Node.js 与 MySQL 的\\[.4cm]在线学习平台设计与实现\par}
\vspace{.8cm}{\Large ——以“知途课堂”为例\par}\vfill
\begin{tabular}{rl}作者：&\underline{\hspace{7cm}}\\[.6cm]学号：&\underline{\hspace{7cm}}\\[.6cm]专业：&\underline{\hspace{7cm}}\\[.6cm]指导教师：&\underline{\hspace{7cm}}\end{tabular}
\vspace{1.5cm}\par 2026年10月8日\end{titlepage}
\frontmatter
\chapter*{摘要}
\addcontentsline{toc}{chapter}{摘要}

针对公开课程入口分散、学习记录难以连续管理的问题，本文设计并实现以“按自己的节奏学习”为主题的知途课堂。平台采用原生HTML、CSS与JavaScript前端、Node.js同源服务和MySQL业务数据库，实现课程发现、账户、收藏、私有笔记、最近访问及每日打卡，并区分第三方视频访问与原生播放器实际进度。本轮实现了课程摘要批量查询、异步密码派生、昵称唯一约束，以及具有服务端权限、事务审计和版本冲突检查的管理工作台。推荐模块包含内容排序和可选BPR模型，采用独立验证与测试分区。2026年10月7日核对的目录包含33门上架课程、2,205个分集；隔离接口与算法测试24项、独立MySQL测试13项、管理端浏览器专项1项通过。33门课程的临时表摘要基准中，查询次数由132次降至4次，中位耗时由28.6毫秒降至1.4毫秒；该结果不代表整站响应时间。业务训练数据不足，合成数据中混合推荐亦未优于内容基线，未启用业务机器学习模型。本文成果支持限定环境内的主要业务流程，不构成学习效果、大规模性能或第三方视频持续可用性的证明。


关键词：在线学习平台；Node.js；MySQL；学习记录；响应式设计


\chapter*{Abstract}
\addcontentsline{toc}{chapter}{Abstract}

This thesis presents Zhitu Classroom, a responsive learning platform built with native web modules, Node.js, and MySQL. It supports course discovery, accounts, bookmarks, private notes, recent visits, and manual daily check-ins, while \path{distinguishing} third-party video visits from native playback progress. The \path{implementation} adds batched course summaries, asynchronous password derivation, a unique nickname constraint, and an \path{administrator} workspace protected by server-side \path{authorization}, \path{transactional} auditing, and version checks. Content \path{recommendations} and an optional BPR model are evaluated using separate validation and test partitions. Verification on October 7, 2026 identified 33 published courses and 2,205 episodes. Isolated API and algorithm tests reported 24 passing records, independent MySQL regression reported 13, and an \path{administrator} browser scenario passed. In a \path{temporary-table} benchmark for 33 course summaries, queries fell from 132 to 4 and median time from 28.6 to 1.4 milliseconds; this is not an end-to-end website latency result. Business data remain insufficient for model training, and hybrid \path{recommendations} did not outperform the content baseline on synthetic data. No business \path{machine-learning} model is enabled. These findings support the tested workflows without establishing learning \path{effectiveness} or \path{production-scale} performance.


Keywords: online learning platform; Node.js; MySQL; learning records; responsive design


\tableofcontents
\mainmatter
\chapter{绪论}

\section{研究背景与问题提出}

在线课程能够为不同学习阶段的用户提供较为灵活的内容选择，但课程资源的增加并不自然等于学习过程的连贯。以跨学科学习为例，用户可能需要在数学课程、编程教程和学习方法文章之间切换。若课程来源、章节入口和笔记分别保存在不同位置，再次进入学习状态时就需要重复查找。对于以公开课程为主要资源的个人学习网站，如何将内容发现与个人记录连接起来，是比单纯增加课程数量更具体的建设问题。


学习支持也不能被简化为播放次数和打卡天数。用户打开某个视频，可能只是查看课程是否符合需要；记录一条笔记，可能发生在视频之外；点击打卡只能说明用户主动完成了这次操作。系统若不区分这些行为，容易将有限的技术信号解释为学习完成或知识掌握。因此，本课题既关注功能实现，也关注学习数据的含义及其展示边界。


本课题围绕三个工程问题展开：如何组织真实来源课程并缩短续学路径；如何在无法取得第三方视频播放秒数的条件下建立合理的学习记录；如何兼顾个人数据持久化、界面便利性和当前项目规模下的可维护性。


\section{相关研究与实践}

自我调节学习为理解学习者主动安排和回顾学习活动提供了研究视角。Zimmerman围绕自我调节学习者展开讨论，其研究可作为本课题分析自主学习行为的理论起点。\cite{ref1} Panadero对六种自我调节学习模型进行了综述，表明该领域涉及认知、元认知、行为及动机等多个方面。\cite{ref2} 对本系统而言，这些研究启发的是为用户保留选择、记录与回顾的空间，而非证明某一种界面组件能够直接提高成绩。本文没有进行学习成效实验，不据此作因果推断。


在平台实践方面，Moodle的课程完成报告允许按照活动表现或人工确认等条件组织完成状态。\cite{ref3} 这说明完成状态需要与明确规则对应。知途课堂使用外部视频来源，无法获得相同的课程活动数据，因而选择将“最近访问”“原生播放进度”和“主动打卡”分开建模。本文不对不同平台作市场排名或性能优劣判断，比较的重点是各自能够取得的数据与功能定位。


现有初稿尚未完成系统性的国内相关文献检索，不能据当前有限资料概括国内外在线学习研究全貌。正式定稿前应补充与学习记录、学习平台设计和可用性评价直接相关的中文研究，并结合学校规定调整综述范围。


\section{研究内容与实施方法}

本课题采用需求梳理、模块化设计、原型迭代和分层验证相结合的方法。首先，根据课程发现、分集学习和再次访问三个连续场景划分功能。其次，建立浏览器、服务端和数据库之间的职责分工，将公开内容、个人业务数据和浏览器偏好分别保存。再次，以不同视频来源的能力差异为依据设计访问、进度与打卡数据。最后，通过隔离接口测试与运行服务检查验证实现，记录无法完成的验证项目。


主要工作包括：构建真实课程目录及分集映射；实现基于会话的账户与个人数据接口；设计学习记录和每日打卡的一致性规则；完成响应式用户页面、搜索与阅读交互；形成数据库结构、维护流程和验证记录。这些是项目工程工作，不等同于提出新的教学理论或推荐算法。


\section{论文结构}

第一章说明研究背景与问题；第二章介绍相关技术及选择依据；第三章分析需求和用例；第四章给出总体架构及接口设计；第五章分析数据库结构；第六章介绍关键功能实现；第七章呈现测试过程与结果边界；第八章总结已完成工作及不足。


\chapter{相关技术与选型依据}

\section{浏览器端模块与响应式布局}

前端采用原生HTML、CSS和JavaScript，使用模块文件组织页面与公共组件。项目规模允许直接由服务器提供资源，不需要额外构建链。该方案有利于追踪DOM更新、路由状态和事件监听器的生命周期，也减少了框架版本变化带来的维护工作。不过，随着页面增加，直接操作DOM会提高状态同步和重复渲染的复杂度，需要通过组件接口和退出清理约定控制影响。


响应式布局使用网格、弹性布局和媒体查询。桌面页面保留较宽的课程列表，手机端采用底部导航和纵向内容排列。动画不是功能成立的前提；系统通过减少动态效果媒体特性识别用户偏好，相应降低动态表现。MDN对该特性的说明是用户已设置减少非必要运动效果，而不是要求所有界面交互失效。\cite{ref4}


\section{Node.js与异步服务}

后端采用Node.js，模块体系为ESM。Node.js官方文档说明了通过import、export及包配置使用ECMAScript模块的方式。\cite{ref5} 本项目在package.json中设置模块类型，通过独立文件划分HTTP服务、业务接口、数据库访问和媒体配置。项目最低版本要求为22.13，本轮实际接口测试使用Node.js v24.16.0。


后端没有把每个页面建设为独立服务，而是采用同源单体应用。该选择适合本地演示和当前数据规模，使页面及接口部署关系直观，但也意味着业务处理错误和资源消耗可能影响同一个进程。数据库查询必须正确等待，密码派生等计算开销也应纳入后续性能评估。


\section{MySQL与数据访问适配}

MySQL作为业务数据库保存关联性较强的数据。用户、课程、章节、收藏和笔记之间存在明确关系，适合通过主键、外键和唯一性约束维护。MySQL文档要求外键两侧固定精度数值字段的大小和符号一致，这与项目旧结构中存在有符号、无符号主键的情况有关。\cite{ref6}


数据库适配层提供异步get、all、run和transaction接口。MySQL侧使用连接池和\path{AsyncLocalStorage}，使事务回调中的查询使用同一连接；SQLite侧用于隔离测试，通过串行请求队列控制访问。该适配层只转换项目使用的少量SQL形式，并不是通用数据库方言转换器。因此，SQLite测试通过后仍需要独立的MySQL测试。


\section{第三方嵌入与同源约束}

浏览器同源策略限制跨来源文档之间的脚本访问；能够取得iframe引用并不意味着能够自由读取其中的播放器状态。\cite{ref7} 当前项目未实现能够取得Bilibili实际播放秒数的原站通信能力。因此，Bilibili课程只记录本站可确认的章节访问，原生直链视频才采用\path{HTMLMediaElement}事件保存播放位置。


同源页面与API减少了部署中的跨域配置需求，但不能替代身份认证、授权或输入校验。本文所说的“同源”是资源组织方式，不是完整安全保证。


\chapter{系统需求分析}

\section{用户角色与业务范围}

系统主要角色包括游客、注册学习者和指定管理员。游客可查看上架课程与公开文章；学习者使用个人学习功能；管理员通过服务端授权管理课程、文章与讨论，维护普通用户状态并查看操作审计。前端页面的显示不能替代服务端权限判断。



\noindent\textbf{表3-1 用户角色与功能边界}

{\small
\begin{longtable}{P{0.32\textwidth}P{0.24\textwidth}P{0.32\textwidth}}
\toprule
角色 & 可用功能 & 限制 \\
\midrule
\endhead
游客 & 课程浏览、公开文章、搜索与原站播放入口 & 私人业务接口要求登录 \\
注册学习者 & 收藏、笔记、互动、头像、历史与打卡 & 不能通过请求参数访问其他用户私有记录 \\
指定管理员 & 数据看板、用户状态、课程编辑、文章发布、讨论审核、审计 & 服务端会话与管理员身份校验，写入MySQL \\
\bottomrule
\end{longtable}}

系统当前不包含在线考试、成绩评定或直播推流。历史直播表与兼容接口保留，页面不再提供直播入口。后台包含课程、文章和讨论三类核心业务管理以及用户管理，不包括任意角色定制、在线授权他人或视频文件托管。


\section{核心功能需求}


\noindent\textbf{表3-2 功能需求与主要验收点}

{\small
\begin{longtable}{P{0.32\textwidth}P{0.24\textwidth}P{0.32\textwidth}}
\toprule
编号 & 功能需求 & 主要验收点 \\
\midrule
\endhead
R01 & 课程发现与检索 & 上架课程可按方向、学科与关键词找到，返回路径明确 \\
R02 & 账户与头像 & 注册登录可恢复会话，头像更新由会话绑定用户 \\
R03 & 章节与来源 & 分集映射明确，来源与播放能力提示准确 \\
R04 & 收藏与笔记 & 收藏可切换，笔记归属与课程章节关系受校验 \\
R05 & 历史记录 & 每门课程显示最近访问，已下架课程不进入常规列表 \\
R06 & 每日打卡 & 服务器确定日期，每账号每天唯一，重复请求不重复计数 \\
R07 & 经验阅读 & 短摘录保留出处，指定全文可在站内阅读和返回 \\
R08 & 响应式与偏好 & 桌面手机布局可用，语言及减少动态效果偏好有响应 \\
\bottomrule
\end{longtable}}

学习过程的典型路径是：在首页或发现页选择课程，进入章节后观看或跳转原站，记录笔记并加入书架；再次访问时从继续学习卡片或个人历史返回。打卡作为独立主动操作，不由视频访问自动触发。搜索历史属于浏览器辅助功能，不能与服务器课程访问历史混为一谈。


\section{非功能需求}

一致性方面，收藏切换、进度更新和打卡写入需要处理并发请求；课程与章节之间还需要应用层联合校验。安全方面，私人接口必须认证，数据库查询使用参数绑定，静态服务不能开放任意项目文件。可维护性方面，页面资源、视频映射和数据库迁移相互分离。易用性方面，弹窗关闭后应恢复焦点，菜单同时支持键盘和触屏，页面退出后应清理计时器和观察器。


本课题未给出“支持万人在线”或固定毫秒级响应等未经测试的指标。正式部署前需要确定硬件条件、数据规模、请求模型和性能目标，再进行负载测试。


\section{可行性分析}

技术上，浏览器媒体、Node.js接口和MySQL关系数据可覆盖当前主要流程；跨域原站视频的进度限制通过缩小记录范围处理。实施上，原生前端与同源服务便于本地运行和逐模块验证。资源上，课程采用有来源的公开目录，服务器主要保存元数据和个人记录，不承担第三方视频转码与分发。由于尚未部署生产服务，本文不估算虚构的年度运营成本，也不把原站公开可访问等同于拥有复制分发权。


\chapter{系统总体设计}

\section{总体架构}

系统采用浏览器、应用服务和数据库三层结构。浏览器端负责界面及本地偏好；应用服务校验会话与业务关系；数据库保存需要跨页面或会话恢复的数据。原站播放器与原生媒体地址是独立内容来源，不经过本站服务器代理视频流。




\begin{figure}[htbp]\centering
\begin{tikzpicture}[box/.style={draw=Line,rounded corners=3pt,fill=Paper,align=center,text width=4cm,minimum height=1cm,font=\small},flow/.style={-{Stealth[length=2mm]},draw=Ink}]
\node[box] (ui) at (0,0) {浏览器页面\\课程、学习、发现、个人中心};
\node[box,text width=3.4cm] (local) at (5,0) {浏览器存储\\偏好、搜索、专注计时};
\node[box] (http) at (0,-1.9) {Node.js HTTP 入口\\同源静态资源与 API};
\node[box,text width=3.4cm] (media) at (5,-1.9) {独立媒体来源\\原站 iframe / 原生直链};
\node[box] (api) at (0,-3.8) {业务层\\认证、关系校验、事务};
\node[box,text width=3.4cm] (db) at (-1.9,-5.7) {MySQL\\账户及业务记录};
\node[box,text width=3.4cm] (file) at (3,-5.7) {头像文件目录};
\draw[flow] (ui)--(http);\draw[flow] (http)--(api);\draw[flow] (ui)--(local);
\draw[flow] (ui.east)--++(.45,0)|-(media.west);
\draw[flow] (api)--(db);\draw[flow] (api)--(file);
\end{tikzpicture}
\caption{系统总体架构}\label{fig:architecture}
\end{figure}


模块边界以变化原因划分。server.mjs负责HTTP入口和白名单，backend.mjs负责业务行为，database.mjs负责异步访问与事务，catalog.mjs负责课程目录规则。页面层将搜索、个人中心、打卡和文章阅读分开，使修改局部交互时无需改写数据库模块。


\section{关键接口设计}

接口通过HTTP方法区分查询和修改。登录后的用户编号从有效会话取得，客户端不能在请求正文中指定笔记或打卡所有者。接口返回JSON及适当状态码，验证失败不应导致整个服务退出。



\noindent\textbf{表4-1 主要接口设计}

{\small
\begin{longtable}{P{0.32\textwidth}P{0.24\textwidth}P{0.32\textwidth}}
\toprule
方法与路径 & 用途 & 身份与关键约束 \\
\midrule
\endhead
GET /api/courses & 课程列表与关键词检索 & 公开读取，只列出上架课程 \\
GET \path{/api/courses/:id} & 课程详情与章节 & 公开信息与当前会话状态分别组织 \\
POST \path{/api/auth/register} & 注册并创建会话 & 校验输入，检查昵称，生成密码摘要 \\
POST \path{/api/me/avatar} & 头像更新 & 要求登录，限制JPEG数据及请求大小 \\
POST \path{/api/courses/:id/visit} & 最近章节访问 & 要求登录，校验课程上架和章节归属 \\
POST \path{/api/courses/:id/progress} & 原生播放进度 & 拒绝Bilibili进度写入，校验时间范围 \\
POST \path{/api/courses/:id/notes} & 创建私有笔记 & 会话绑定用户，校验章节及内容 \\
DELETE \path{/api/notes/:id} & 删除自己的笔记 & 删除条件同时包含记录与当前用户编号 \\
GET、POST \path{/api/me/checkins} & 查询与提交打卡 & POST只接受空对象，日期和账号由服务端确定 \\
\bottomrule
\end{longtable}}

已下架课程不再展示在常规目录中，但详情兼容行为不等同于统一封禁所有旧接口。现状下，访问写入明确检查上架状态，其他旧详情和互动路径仍需在正式运营前统一审视。论文不将保留兼容接口描述为完整的内容访问控制体系。


\section{数据归属设计}

MySQL保存身份、收藏、笔记、访问、打卡、管理员身份、文章、评论审核及审计记录。头像文件单独保存，users.avatar只记录访问地址。搜索历史与界面偏好保存在浏览器，专注计时使用标签页会话存储。九篇精选及《提问的智慧》仍为静态资源；后台新增的学习随笔保存到MySQL，发布后在发现页展示。


本地搜索历史和偏好不会自动跨浏览器迁移，专注计时也不能视为真实课程观看时长。管理数据经服务端写入，与这些浏览器偏好具有不同持久化边界。


\section{一致性与异常处理}

MySQL事务从连接池获取连接，启动事务后通过异步上下文传递给查询方法，在回调完成时提交，出现异常时回滚并释放连接。涉及同一用户状态变更时，使用用户行锁顺序处理。在当前项目中，该方法便于统一收藏切换和打卡的并发保护；代价是同一用户不同功能的写入也可能等待同一把锁。


数据库结构修改不与普通业务事务混同。MySQL文档列举了会导致隐式提交的语句，包括多种DDL操作。\cite{ref8} 因此，结构升级需要单独备份和验证，不能依靠业务事务回滚撤销已经执行的全部建表或改表操作。


\chapter{数据库设计}

\section{概念模型}

核心实体为用户、课程、章节和笔记。用户与课程之间通过收藏、访问和进度等关联记录建立不同含义的关系；每日打卡直接关联用户与业务日期。图5-1展示当前学习流程中的主要关系，不包含历史直播、会话和所有互动表。




\begin{figure}[htbp]\centering
\begin{tikzpicture}[box/.style={draw=Line,rounded corners=2pt,fill=Paper,align=center,text width=2.35cm,minimum height=.85cm,font=\small},rel/.style={-{Stealth[length=1.7mm]},draw=Ink},every edge quotes/.style={font=\scriptsize,fill=white,inner sep=1pt}]
\node[box] (u) at (0,0) {用户};\node[box] (c) at (8,0) {课程};
\node[box] (f) at (4,0) {收藏};\node[box] (ch) at (8,-2) {章节};
\node[box] (v) at (4,-2) {最近访问};\node[box] (n) at (4,-4) {私有笔记};
\node[box] (p) at (8,-6) {章节播放};\node[box] (d) at (0,-6) {每日打卡};
\draw[rel] (u)--node[above,font=\scriptsize]{1:N}(f);
\draw[rel] (c)--node[above,font=\scriptsize]{1:N}(f);
\draw[rel] (c)--node[right,font=\scriptsize]{1:N}(ch);
\draw[rel] (u.south)--++(0,-1)|-(v.west);
\draw[rel] (u.south)--++(0,-4)--(n.west);
\draw[rel] (u)--node[left,font=\scriptsize]{1:N}(d);
\draw[rel] (u.west)--++(-.5,0)|-([yshift=-.9cm]p.south)--(p.south);
\draw[rel] (c.east)--++(.65,0)|-(v.north east);
\draw[rel] (c.east)--++(1,0)|-(n.south east);
\draw[rel] (ch)--(v);\draw[rel] (ch)--(n);\draw[rel] (ch)--(p);
\end{tikzpicture}
\caption{核心学习数据关系（局部，箭头从父表指向关联记录，均为一对多）}\label{fig:relations}
\end{figure}


以收藏为例，关系键为用户与课程；以最近访问为例，用户与课程组合只对应一条最近记录，而不是无限追加日志。不同关系不能仅因为都包含用户编号和课程编号就合并为同一种记录，否则会混淆收藏意愿、打开行为与播放位置。


\section{物理结构与关键字段}

2026年10月7日采集的实际MySQL结构包含22张基础表、102个基础表字段、30个外键和1个视图。\path{history_resets}阻止清空历史后由旧progress补出访问；本轮新增\path{admin_accounts}、\path{account_status}、\path{admin_articles}、\path{comment_moderation}和admin\_audit五张表，并补建昵称唯一索引。升级前检查重复昵称及父列类型，备份相关结构；未覆盖既有学习记录。



\noindent\textbf{表5-1 主要数据表}

{\small
\begin{longtable}{P{0.32\textwidth}P{0.24\textwidth}P{0.32\textwidth}}
\toprule
表名 & 主键 & 核心作用 \\
\midrule
\endhead
users & id & 用户昵称、密码摘要、头像地址 \\
courses & id & 标题、分类、学科、讲师、上架状态 \\
chapters & id & 所属课程、章节标题、顺序、时长 \\
favorites & user\_id + course\_id & 用户课程收藏关系 \\
notes & id & 私有笔记、所属章节及时间点 \\
\path{course_visits} & user\_id + course\_id & 每门课程最近打开的章节 \\
\path{chapter_playback} & user\_id + chapter\_id & 原生视频位置及章节进度 \\
\path{study_checkins} & user\_id + checkin\_date & 每用户每天唯一的主动打卡 \\
\path{history_resets} & user\_id & 清除历史的时间标记 \\
\path{admin_accounts} & user\_id & 指定管理员身份 \\
\path{account_status} & user\_id & 普通账户停用状态 \\
\path{admin_articles} & id & 纯文本文章与发布状态 \\
\path{comment_moderation} & comment\_id & 评论可见状态 \\
admin\_audit & id & 操作者、动作、对象与时间 \\
sessions & token & 用户会话和失效时间 \\
\bottomrule
\end{longtable}}

其余基础表为comments、danmaku、follows、coins、progress、live\_events、reservations，视图为\path{v_course_comment_stats}。直播相关表仅为历史兼容。讲师关注当前以讲师名称建立关系，没有独立讲师账户实体；弹幕表也没有用户编号列，不能将其描述为完整的用户行为审计表。



\noindent\textbf{表5-2 笔记与打卡关键字段摘录}

{\small
\begin{longtable}{P{0.32\textwidth}P{0.24\textwidth}P{0.32\textwidth}}
\toprule
表与字段 & 类型 & 约束或含义 \\
\midrule
\endhead
notes.id & BIGINT & 自增主键，有符号 \\
\path{notes.user_id} & BIGINT UNSIGNED & 关联用户，私人归属 \\
\path{notes.course_id} & BIGINT UNSIGNED & 关联课程 \\
\path{notes.chapter_id} & BIGINT UNSIGNED & 关联章节，仍需应用层校验课程匹配 \\
\path{notes.at_second} & DOUBLE & 视频时间点，单位为秒 \\
\path{notes.content} & TEXT & 笔记正文，接口限制长度 \\
\path{study_checkins.user_id} & BIGINT UNSIGNED & 联合主键组成部分，关联用户 \\
\path{study_checkins.checkin_date} & DATE & 联合主键组成部分，北京时间业务日期 \\
\path{study_checkins.created_at} & DATETIME(3) & 写入时间，不代替业务日期 \\
\bottomrule
\end{longtable}}

表中字段类型取自实际结构快照。接口内容长度与时间范围检查不等同于数据库CHECK约束。\path{users.nickname}已具有单列唯一索引，注册前查询用于友好提示，最终由数据库拒绝并发重复插入；同名判断采用数据库排序规则。


\section{三类学习记录的分离}

第一类记录是\path{course_visits}，表示用户最近打开了哪门课程的哪个章节。第二类是\path{chapter_playback}及progress，保存原生播放器可报告的播放位置与汇总进度。第三类是\path{study_checkins}，表示某个业务日期进行了主动打卡。三个模型对应不同数据来源与语义，不能相互代替。


原生章节进度按当前播放位置和时长计算，写入时保留该章节已记录的最大百分比，但最近位置可随拖动而变化。课程进度由所有章节的已记录百分比平均得到，未记录章节按零计算。它属于媒体位置指标，并不证明用户认真观看或掌握内容，也不是课程考试成绩。对于混合来源课程，原站章节不能自动回传进度，这会影响汇总指标的解释。


\section{数据维护策略}

业务数据默认来自MySQL，历史SQLite文件不再代表最新状态。课程目录调整使用限定编号的增量同步，先预检和备份再应用；明确下架采用published字段，不直接删除关联章节。头像备份同时包含数据库地址和文件目录。修改页面、颜色和文档时，不运行全量迁移。


\chapter{关键功能实现}

\section{认证与头像更新}

注册先校验输入并查询昵称占用，再用随机盐与异步scrypt派生摘要。\cite{ref9} 与既有摘要格式保持兼容。派生队列最多运行4个任务、等待16个，超限返回429；登录及注册对直接连接地址限制每15分钟60次。该限流是单进程内存状态，不支持分布式协调，反向代理场景需另行设计可信地址处理。8个并发同名注册在隔离SQLite与MySQL用例中均只有一次成功，其余返回409；测试不代表所有并发故障都已覆盖。


登录验证成功后，服务端生成随机会话令牌并写入sessions，Cookie设置HttpOnly、SameSite=Lax和失效时间。请求处理中先读取有效会话，再确定用户编号。当前是本地HTTP运行方式，不能把该Cookie设置描述为已经完成HTTPS部署或所有CSRF风险防护。


头像由浏览器裁剪后提交JPEG数据，服务端限制请求和内容形式，生成随机文件名并更新当前用户头像地址。若数据库写入失败，则清理本次新建文件，避免产生无关联文件。服务只允许通过限定格式的头像接口读取文件，不开放上传目录浏览。


\section{搜索建议与历史管理}

搜索组件包含空闲轮播和输入检索两种状态。空闲状态展示真实课程名称，用户点击搜索按钮后可直达当前课程；输入状态取消轮播，等待短暂输入停顿后向接口请求建议。代码通过约180毫秒的延迟减少连续键入造成的重复请求，并使用\path{AbortController}和递增请求编号丢弃过时结果。


搜索历史按当前账号或游客身份形成不同本地存储键，清空操作仅删除对应历史。垃圾桶按钮保留文字形式的title和aria-label，既满足紧凑排版，也保留操作含义。课程页的大搜索框被移除后，顶部成为统一入口，降低多个输入状态不同步的风险。


当前检索由SQL模糊匹配实现，并非全文搜索引擎。检索与推荐分别实现：首页推荐已改为基于内容的个性化排序，并提供可离线训练的BPR矩阵分解模块。课程规模进一步增长时，需要测量查询开销并考虑分页、索引或全文检索。


\section{推荐方法与查询优化}

内容推荐将课程类别、学科、讲师以及标题和简介词项编码为TF-IDF向量，中文采用双字切分。收藏、投币和最近访问分别赋予3、2、1的工程权重，同一用户—课程信号合并后封顶为6；对课程向量进行加权聚合得到用户兴趣向量。推荐分数由余弦相似度与归一化关注人数组成，权重分别为0.9和0.1。优先展示未接触的课程，只返回已上架课程。无行为用户依据去重关注人数排序，无关注时按编号稳定排序。这些参数是当前实现选择，尚未经真实用户实验调优。


机器学习模块采用BPR隐反馈排序思想\cite{ref10}，以用户与课程的8维隐向量内积表示潜在偏好，使用随机梯度更新优化已观察课程相对于未观察课程的排序，并加入正则项。未观察项仅用于相对排序采样，不等于用户明确负面反馈。混合分数由内容相似度、经过sigmoid变换的隐向量内积以及关注人数组成，权重分别为0.65、0.25和0.1。该实现属于已有方法的工程应用，不宣称提出新的推荐算法。


训练命令npm run \path{recommend:train}只读收藏、投币和访问数据，不使用私有笔记正文。至少10名有行为用户、50组去重交互和3门课程才进入评估；将去重记录按全局唯一时间戳分组，以约60\%、20\%、20\%划分训练、验证、测试，边界不拆分同一时间戳；每分区至少5名用户。训练集拟合后，在验证集比较热门、内容、纯BPR与混合四种方法。仅混合NDCG@5严格超过热门与内容基线且HitRate@5不低于两者，才通过启用门控。随后用训练加验证记录拟合并仅评估一次测试集，不用测试指标反复选择参数；通过门控时另用全量记录拟合供服务使用。固定8维、160轮、种子42与0.25模型权重，未做超参数搜索。


2026年10月7日业务训练预检为3位有行为用户、15组去重交互，未达到门槛，当前使用内容推荐。模型与数据库范围绑定、私有保存，重启载入；过期7天、格式无效或用户无可用向量时回退。模型文件及用户隐向量不公开，也不随论文源码附件分发。


这一全局分区仍基于“最近用户—课程状态快照”：收藏取消与早期访问已不可见，元数据没有历史版本，未建立曝光日志。因此消除了记录时间戳的跨分区倒置，并不等同完整事件流的严格历史回放，也不能得出线上因果收益。指标按有未见相关课程的可评估用户宏平均；Coverage@5为推荐并集占上架目录比例。合成实验与业务预检分别报告于第七章。


课程摘要原先对每门课程分别查询进度、最近章节、收藏和投币，共4N次。本轮改为按课程编号集合分别执行四类参数绑定查询，再以映射表拼装结果，空目录或匿名用户无需查询个人状态。课程内容向量按元数据签名缓存，内容变化后重算；兴趣状态仍按请求读取，避免跨用户复用。批处理不改变私有数据的用户条件，也不等同数据库全文搜索。


\section{视频访问与笔记归属}

原站视频由bvid、cid和分P信息组成映射。页面加载某章节时选择正确的视频来源，并向访问接口提交章节编号。服务端检查课程是否上架以及章节是否属于该课程，成功后覆盖该用户在该课程上的最近章节。历史记录展示最近访问日期，避免用“已完成”文案误导用户。


笔记接口在检查身份后校验内容、章节与时间点。删除采用同时匹配笔记编号和用户编号的SQL条件，避免只依据前端显示隐藏按钮来保护数据。私有笔记查询同样在服务端限定用户；Markdown导出只导出当前页面筛选得到的个人笔记，并提供返回对应章节的链接。


\section{每日打卡的幂等处理}

每日打卡以服务器按UTC加8小时计算的日期为准，请求正文不能传入账号或日期。服务端在事务中锁定当前用户，查询当天是否已经存在记录，再决定是否插入；联合主键作为数据库约束，防止同用户同日期出现重复行。重复提交返回同一天状态，表现为幂等，而收藏切换是另一种交互语义，重复请求可能切换回原状态。




\begin{figure}[htbp]\centering
\begin{tikzpicture}[box/.style={draw=Line,rounded corners=3pt,fill=Paper,align=center,text width=5.2cm,minimum height=.85cm,font=\small},decision/.style={diamond,aspect=3,draw=Line,align=center,font=\small,inner sep=2pt},flow/.style={-{Stealth[length=2mm]},draw=Ink}]
\node[box] (a) at (0,0) {收到打卡请求};
\node[decision] (b) at (0,-1.5) {有效会话且正文为空？};
\node[box,text width=2cm] (x) at (5,-1.5) {返回错误};
\node[box] (c) at (0,-3) {服务器计算北京时间日期};
\node[box] (d) at (0,-4.4) {开启事务并锁定用户行};
\node[decision] (e) at (0,-5.9) {当天记录已存在？};
\node[box,text width=2cm] (f) at (5,-5.9) {插入记录};
\node[box] (h) at (0,-7.5) {提交，查询近 365 天记录};
\node[box] (i) at (0,-8.9) {返回日历与统计};
\draw[flow] (a)--(b);\draw[flow] (b)--node[above,font=\small]{否}(x);
\draw[flow] (b)--node[left,font=\small]{是}(c);\draw[flow] (c)--(d);\draw[flow] (d)--(e);
\draw[flow] (e)--node[above,font=\small]{否}(f);\draw[flow] (e)--node[left,font=\small]{是}(h);
\draw[flow] (f)|-(h);\draw[flow] (h)--(i);
\end{tikzpicture}
\caption{每日打卡处理流程}\label{fig:checkin}
\end{figure}


累计天数使用滚动窗口内不同打卡日期的数量，连续天数从今天开始回溯；若今天尚未打卡，则从昨天开始回溯，直到遇到空缺。所有统计只覆盖近365天，不能称为终身累计。客户端以日期网格展示结果，颜色区分是否打卡，不以颜色深浅暗示学习强度。


\section{响应式界面与阅读模块}

平台采用浅灰背景、白色卡片、粉色点缀和深蓝精选横幅。图形标志由圆润Z字路线和两枚节点组成，强调沿途积累。首页逐字书写、导航滑块和背景动效围绕主题组织，并在减少动态效果偏好下调整。课程卡片和个人历史分别采用适合发现与回访的布局；历史封面保持16:9，卡片沿水平时间轴排列，支持触屏、方向按钮与键盘；时间轴显示北京时间的最近访问时刻。清除按钮需确认，并在当前用户事务中删除访问记录、保存清除标记，保留收藏、笔记与播放进度。


经验区直接展示九篇精选，保留简短英文摘录、中文译文和出处。《提问的智慧》则使用用户明确提供的完整文件，静态生成正文并校验内部锚点。全文以dialog呈现，目录点击只滚动文章，不改变站点路由；关闭后恢复入口焦点。原文内容、站内界面翻译和用户笔记分别处理，避免语言切换改写引用文本。


前端路线切换还涉及生命周期管理。学习页、个人工具中的观察器与计时器需要随离开页面而清理；页头搜索为跨路由常驻组件，生命周期与普通页面不同。将这一差异明确化，有助于减少重复监听和后台计时问题。


\section{管理授权与发布流程}

后台以\path{admin_accounts}记录受明确指定的管理员，注册请求不能自行传入角色。管理接口每次检查有效会话与管理员身份，匿名返回401，普通账号返回403。当前提供数据看板、用户状态维护、课程编辑与上下架、文章编写发布、评论隐藏恢复及审计查询，均使用真实业务API。授权由运维脚本按唯一昵称配置，不提供浏览器内自我授权。


编辑接口在事务中锁定操作者与目标行，以返回给页面的行摘要版本核对当前值；陈旧提交返回409，避免两个编辑者静默覆盖。变更与审计同行事务提交，审计保存操作者、动作、目标和时间，不记录密码或全文副本。停用普通账号会撤销其会话，管理员账号不能由该界面停用。课程只调整元数据与上架状态，不删除章节或修改片源；已明确退役的课程不能直接重新发布。文章正文以纯文本转义显示；草稿接口不可公开读取，发布后在发现页展示，取消发布后不再公开。评论隐藏保留原记录。


\chapter{系统测试与结果分析}

\section{验证环境与范围}

验证日期为2026年10月7日，环境为Windows x64、Node.js v24.16.0、MySQL 8.0.30。接口与浏览器使用临时SQLite、隔离头像目录和测试账号；MySQL回归启动独立临时实例。业务库仅用于结构检查、授权升级及训练汇总，不把真实用户笔记、头像或会话用于用例。当前业务目录为33门上架课程、2,205个分集。



\noindent\textbf{表7-1 统一验证口径}

{\small
\begin{longtable}{P{0.32\textwidth}P{0.24\textwidth}P{0.32\textwidth}}
\toprule
命令或检查 & 本轮结果 & 适用边界 \\
\midrule
\endhead
npm test & 24项通过，0失败、0跳过 & 临时SQLite接口与算法，包括运行器父测试 \\
npm run test:mysql & 13项通过，0失败、0跳过 & 独立临时MySQL，含父测试，不是业务库压力测试 \\
ADMIN\_ONLY=1浏览器专项 & 1项通过 & Edge，后台桌面/手机与公开文章阅读 \\
db:check及结构采集 & MySQL可访问，22表、1视图 & 不含业务数据的实际元数据 \\
全量浏览器回归 & 本轮未重跑 & 专项通过不能替代全站交互回归 \\
HTTP并发与学习效果 & 未开展 & 不作生产性能或学习收益结论 \\
\bottomrule
\end{longtable}}

此前历史时间轴专项已通过，证据保存在\path{evidence/history-validation.md}；早期19/11及20/12等计数属于不同代码版本，本章统一使用本轮24/13结果。运行器父测试计入上述数字，但不能当作额外独立场景。完整证据索引为\path{evidence/platform-validation.md}。


\section{业务一致性与权限验证}

基础回归覆盖静态资源白名单、账户及头像、课程分集、私有笔记、进度与访问区别、每日打卡、历史清除和持久化。新增用例中，8个并发同名注册仅1个成功，7个返回409；注册提交管理员字段不能获得权限。匿名后台访问返回401，普通用户返回403。后台课程修改可持久化，旧版本再次提交返回409，退役课程拒绝重新上架。文章草稿不可公开读取，发布后可读取，取消发布后返回404，非HTTP(S)来源地址被拒绝。


评论隐藏及恢复影响公开列表；停用账号撤销会话且拒绝重新登录，恢复后可重新登录；管理员不允许通过该界面停用。管理员用户列表不含密码摘要或令牌，成功写入产生审计记录。上述用例在SQLite与独立MySQL均执行。浏览器专项覆盖后台课程编辑、文章发布及移动阅读，正文中的脚本文字按文本显示。测试没有模拟所有数据库宕机或跨进程竞态。


\section{摘要查询优化实验}

基准命令为npm run \path{benchmark:summaries}，使用配置的MySQL连接，但创建的均为连接私有TEMPORARY表和合成记录，不读取或改写业务账户。旧方法逐课程执行四类查询，新方法一次批量执行四类查询，核对返回数量、进度与章节标题一致。每组预热3次、交替顺序测量20次，报告中位数和经验P95。原始结果见\path{evidence/summary-benchmark.json}。



\noindent\textbf{表7-2 单连接课程摘要基准}

{\small
\begin{longtable}{P{0.176\textwidth}P{0.176\textwidth}P{0.176\textwidth}P{0.176\textwidth}P{0.176\textwidth}}
\toprule
合成课程数 & 方法 & SQL次数 & 中位耗时/ms & P95/ms \\
\midrule
\endhead
33 & 逐课程 & 132 & 28.576 & 29.331 \\
33 & 批量 & 4 & 1.377 & 1.642 \\
330 & 逐课程 & 1320 & 300.742 & 307.780 \\
330 & 批量 & 4 & 6.421 & 6.772 \\
\bottomrule
\end{longtable}}

本实验支持减少摘要函数查询往返的结论。它使用同一连接与构造数据，没有模拟业务连接池、网络抖动、HTTP并发、客户端渲染或真实索引分布，因此不能将该比例写为整站加速比例或部署容量。样本仅20次，分位数稳定性有限；真实负载评价需要独立请求模型和更长测量。


\section{推荐对照实验}

npm run \path{recommend:experiment}使用40用户、40课程、320组交互的合成数据，按唯一时间戳得到160/80/80条训练、验证、测试记录。所有方法采用相同可见数据、未见上架候选与K=5。该构造具有明确主题聚类，便于重复验证程序与门控，容易取得高分，不代表实际学习者分布。原始结果见\path{evidence/recommendation-synthetic.json}。



\noindent\textbf{表7-3 独立测试分区结果（合成样本）}

{\small
\begin{longtable}{P{0.176\textwidth}P{0.176\textwidth}P{0.176\textwidth}P{0.176\textwidth}P{0.176\textwidth}}
\toprule
方法 & HitRate@5 & Recall@5 & NDCG@5 & Coverage@5 \\
\midrule
\endhead
热门 & 0.250 & 0.250 & 0.229 & 0.275 \\
内容 & 1.000 & 1.000 & 0.914 & 1.000 \\
纯BPR & 1.000 & 1.000 & 0.774 & 1.000 \\
混合 & 1.000 & 1.000 & 0.774 & 1.000 \\
\bottomrule
\end{longtable}}

验证集内容NDCG@5为0.762，混合为0.728；混合未优于内容，门控结果为\path{baseline-not-improved}，未生成启用模型。测试集也没有给出混合优于内容的证据，不应为了获得正向结果反复调参。业务数据预检仅3用户、15组交互，返回\path{insufficient-data}，业务效果指标尚不存在。本实验没有多次随机种子置信区间、曝光日志或在线A/B结果。


\section{有效性限制与后续评价}

结果支持限定样本中的数据归属、批量查询语义、并发唯一注册、权限与管理写入。完整浏览器回归本轮未重跑，专项通过不能替代所有页面、辅助技术及浏览器版本的验证。第三方视频是否持续可播放仍受原站影响。下一步应补充真实负载的HTTP分位数与错误率、完整浏览器流程、备份恢复演练，并在知情同意下积累事件和曝光记录。真实用户任务评价及学习效果研究需要另行设计，不预设满意度或成绩改善。


\chapter{总结与展望}

\section{工作总结}

本文围绕公开课程的发现与持续访问需求，实现了基于Node.js和MySQL的知途课堂。系统将课程目录、章节来源、账户、收藏、私有笔记、历史记录、主动打卡和学习经验组织在同一响应式网站内，形成从寻找课程到回顾学习记录的主要流程。


项目对数据含义和权限边界作了区分：访问不等于播放完成，打卡不等于学习强度，个人工具状态与正式管理写入分开。批量摘要、异步密码派生和昵称唯一约束改善实现基础；服务端管理授权、版本检查和事务审计形成课程、文章、讨论维护流程。内容推荐与可选BPR均有可运行实现，但业务数据不足且合成实验未显示混合优势，故保留内容推荐。


\section{不足与后续工作}

系统仍需完善四方面工作：第一，管理权限仅有学习者与指定管理员两类，尚无细分运营角色、授权撤销界面或完整故障恢复演练。第二，认证限流和派生队列为单进程设计，尚无分布式负载与资源饱和实验。第三，推荐数据为最近关系快照，缺乏完整事件和曝光记录，真实效果未知。第四，课程外部可用性、全量浏览器回归及真实学习者体验还需持续评价。


后续工作应先保证数据和权限正确，再扩展内容运营与部署。正式上线还需补充HTTPS、代理环境限流、备份恢复演练和监控。学习效果研究属于另一层评价，只有在建立明确研究问题、样本与评价设计后，才能讨论系统对学习行为或成绩的影响。


\begin{thebibliography}{99}
\addcontentsline{toc}{chapter}{参考文献}
\bibitem{ref1} ZIMMERMAN B J. Becoming a \path{Self-Regulated} Learner: An Overview[J]. Theory Into Practice, 2002, 41(2): 64-70. DOI: \href{https://doi.org/10.1207/s15430421tip4102_2}{10.1207/s15430421tip4102\_2}.
\bibitem{ref2} PANADERO E. A Review of \path{Self-regulated} Learning: Six Models and Four Directions for Research[J]. Frontiers in Psychology, 2017, 8: 422. DOI: \href{https://doi.org/10.3389/fpsyg.2017.00422}{10.3389/fpsyg.2017.00422}.
\bibitem{ref3} Moodle. Course completion[EB/OL]. [2026-10-05]. \href{https://docs.moodle.org/en/Course_completion}{Moodle官方文档}.
\bibitem{ref4} MDN \path{contributors.} \path{prefers-reduced-motion}[EB/OL]. [2026-10-05]. \href{https://developer.mozilla.org/en-US/docs/Web/CSS/Reference/At-rules/@media/prefers-reduced-motion}{MDN文档}.
\bibitem{ref5} Node.js. Modules: ECMAScript modules, Node.js v22 \path{documentation}[EB/OL]. [2026-10-05]. \href{https://nodejs.org/docs/latest-v22.x/api/esm.html}{Node.js官方文档}.
\bibitem{ref6} Oracle. MySQL 8.0 Reference Manual: FOREIGN KEY Constraints[EB/OL]. [2026-10-05]. \href{https://dev.mysql.com/doc/refman/8.0/en/create-table-foreign-keys.html}{MySQL官方文档}.
\bibitem{ref7} MDN \path{contributors.} Same-origin policy[EB/OL]. [2026-10-05]. \href{https://developer.mozilla.org/en-US/docs/Web/Security/Defenses/Same-origin_policy}{MDN文档}.
\bibitem{ref8} Oracle. MySQL 8.0 Reference Manual: Statements That Cause an Implicit Commit[EB/OL]. [2026-10-05]. \href{https://dev.mysql.com/doc/refman/8.0/en/implicit-commit.html}{MySQL官方文档}.
\bibitem{ref9} Node.js. Crypto: \path{crypto.scrypt}(), Node.js v22 \path{documentation}[EB/OL]. [2026-10-07]. \href{https://nodejs.org/docs/latest-v22.x/api/crypto.html\#cryptoscryptpassword-salt-keylen-options-callback}{Node.js官方文档}.
\bibitem{ref10} RENDLE S, \path{FREUDENTHALER} C, GANTNER Z, et al. BPR: Bayesian Personalized Ranking from Implicit Feedback[C]. Proceedings of the Twenty-Fifth Conference on Uncertainty in Artificial Intelligence, 2009: 452-461. \href{https://arxiv.org/abs/1205.2618}{论文公开版本}.
\end{thebibliography}
\chapter*{致谢}
\addcontentsline{toc}{chapter}{致谢}

本节待作者根据实际指导、协作和支持情况撰写。正式提交前补充真实姓名与贡献，不使用虚构经历。


\appendix
\chapter{源码与论文论述对应}

\section{A.1 项目源码下载}

\href{zhitu-classroom-source.zip}{下载知途课堂项目源码（ZIP）}


本论文附带2026年10月8日的源码快照，压缩包文件名为\path{zhitu-classroom-source.zip}。将该文件与论文置于同一目录，可通过上述相对链接访问；若阅读器不支持附件链接，可按文件名直接解压。此链接为随论文提交的附件路径，不是公开互联网仓库地址。


源码包包含前后端源程序、课程公开元数据、页面资源、数据库结构、测试脚本与启动说明，不包含实际业务数据库、个人上传头像、密码配置和私人备份。具体文件与SHA-256值见压缩包内\path{SOURCE-MANIFEST.json}。复现请先阅读包内\path{README-SOURCE.md}；测试结果的适用范围仍以第七章为准。


\section{A.2 论述与实现对应}

{\small
\begin{longtable}{P{0.32\textwidth}P{0.24\textwidth}P{0.32\textwidth}}
\toprule
论文主题 & 项目文件 & 核对重点 \\
\midrule
\endhead
服务与认证 & server.mjs、backend.mjs、\path{passwords.mjs} & 白名单、会话与私人数据归属 \\
管理与文章 & \path{admin-api.mjs}、admin.js、\path{published-articles.js} & 服务端角色、版本冲突、发布与审计 \\
推荐与评估 & \path{recommendations.mjs}、\path{recommendation-evaluation.mjs} & 四算法对照与独立分区 \\
批量摘要 & \path{course-summary.mjs} & 查询次数与用户条件 \\
数据库与事务 & database.mjs、\path{mysql-schema.mjs} & 异步连接、事务、实际结构 \\
课程与视频 & catalog.mjs、media.mjs、player.js & 上架、分集映射及进度来源 \\
打卡 & checkins.mjs、\path{checkin-calendar.js} & 北京时间、滚动窗口与展示 \\
个人学习 & profile.js、\path{study-collection.js} & 书架、历史与笔记导出 \\
搜索与阅读 & search.js、\path{study-posts.js} & 轮播、历史、文章与来源 \\
实际数据字典 & database/数据库表结构说明.md & 字段、外键、索引和采集日期 \\
本轮测试证据 & \path{thesis/evidence/} & 测试输出、数据库检查、环境说明 \\
\bottomrule
\end{longtable}}

\chapter{定稿前补充事项}

\begin{enumerate}
\item 根据学校模板填写封面、声明、格式和实际作者信息，确认字数与引用规范。
\item 补充中文研究文献，扩展文献综述；阅读原文后再提出与本课题相关的结论。
\item 使用隔离演示账号截取首页、课程、学习、个人中心与全文阅读页面，隐藏私人记录。
\item 对图4-1、图5-1和图6-1按最终实现校对，局部数据关系图不替代完整物理ER图。
\item 保持MySQL回归，补足全量浏览器与HTTP负载实验；现有摘要基准和合成推荐实验只支持各自限定结论。
\item 核实项目分工与独立完成范围，按学校要求说明使用的工具和辅助方式。
\item 当前文件为初步草稿，不能将待补工作、演示功能或技术假设写成已验证成果。
\end{enumerate}

\end{document}
